% Tutorial solution writeup — the full worked solution set for one sheet. Lives in that sheet's
% folder under 20 Tutorials/, and a past paper takes the same shape under 40 Past Papers/.
% Copy it there and replace every <...> placeholder; the input line below needs no edit.
\documentclass[a4paper,10pt]{article}
\IfFileExists{preamble.tex}{\input{preamble.tex}}{\input{../../templates/preamble.tex}}
\begin{document}
\ArtifactHeader{\ModuleCode}{<Source-map key>}{<the sheet's own name> --- Solutions}{<YYYY-MM-DD>}%
  {<the sheet, and whatever verified solutions this was checked against>}

\section{<The sheet's own grouping, or Solutions where it has none>}

% The optional argument takes the sheet's own numbering, which is what a reader looks for.
\begin{question}[<the sheet's number>]
<The question, restated closely enough that the solution reads without the sheet in hand.>
\end{question}

\begin{solution}
<The working, in the steps a reader would have to produce themselves — the step that is easy to
skip is the step worth writing.>
\[
  \text{<the display the answer turns on>}
\]
<The answer, stated as an answer.>
\end{solution}

\begin{question}[<the sheet's number>]
<A question whose difficulty is a trap rather than the algebra.>
\end{question}

\begin{solution}
<The working.>
\end{solution}

\begin{warning}
<The trap that question sets, and how to see it coming next time. This belongs here only when it
is about the question; what the sheet taught goes in the concepts consolidation.>
\end{warning}
\end{document}
